% Generated by `taskgen course examples` from factoring.pm1_backtrack (seed 0).
% Do not edit: change examples.yaml and regenerate.
\begin{example}[алгоритм $p - 1$, случай $\gcd = N$]\label{ex:pm1-backtrack}
Число $N = 2501$ раскладывается алгоритмом $p - 1$ Полларда с основанием $a = 2$ и границей $B = 10$, т.\,е. с показателем $M = \operatorname{lcm}(1, \ldots, 10) = 2^{3} \cdot 3^{2} \cdot 5 \cdot 7 = 2520$. Вычисляя $a^M$ последовательным возведением в степени $\ell^{e_\ell}$, показать, что $\gcd(a^M - 1, N) = N$, и всё же найти нетривиальный делитель $N$ (возвратом к промежуточным значениям; при неудаче --- следующее основание $a = 3, 5, \ldots$).

\par\noindent\textit{Решение.} $x_0 = a$, $x_i = x_{i-1}^{\ell_i^{e_i}} \bmod N$ по простым $\ell_i \le B$, $g_i = \gcd(x_i - 1, N)$. Если $p \mid N$ и $\operatorname{ord}_p a \mid \ell_1^{e_1} \cdots \ell_i^{e_i}$, то $p \mid g_i$; $g = N$ означает, что порядки $a$ по модулю обоих простых делителей разделили показатель одновременно.\par\noindent $a = 2$:\[\begin{array}{c|rrrr}\ell^{e} & 8 & 9 & 5 & 7 \\ \hline x_i & 256 & 1595 & 1 & 1 \\ g_i & 1 & 1 & 2501 & 2501\end{array}\]$\gcd(a^M - 1, N) = N$. Первое $g_i > 1$ --- при $\ell^e = 5$: $g = 2501$, т.\,е. $g = N$. Возвращаемся к $x = 1595$ и возводим в степень $5$ по одному разу: $x^{5^{1}} \equiv 1$, $\gcd(1 - 1, N) = 2501$. Снова $N$: при этом основании оба порядка дописываются одним и тем же множителем $5$ --- меняем основание.\par\noindent $a = 3$:\[\begin{array}{c|rrrr}\ell^{e} & 8 & 9 & 5 & 7 \\ \hline x_i & 1559 & 985 & 1 & 1 \\ g_i & 41 & 41 & 2501 & 2501\end{array}\]$\gcd(a^M - 1, N) = N$. Первое $g_i > 1$ --- при $\ell^e = 8$: $g = 41$, это и есть делитель (возврат к промежуточному значению).\par $N / 41 = 61$; действительно, $41 - 1 = 2^{3} \cdot 5$ и $61 - 1 = 2^{2} \cdot 3 \cdot 5$ --- оба $10$-гладкие (делят $M$).\[N = 41 \cdot 61 = 2501.\]
\end{example}
